% ------------------------------------------------------------------------------
%						2 Job Management
% ------------------------------------------------------------------------------
%
We use SLURM as the workload manager.
It supports primarily two types of jobs: batch and interactive.
Batch jobs are used to run unattended tasks, whereas, 
interactive jobs are are ideal for setting up virtual environments, compilation, and debugging.\\
%
\noindent \textbf{Note:} In the following instructions, anything bracketed like, \verb+<>+, indicates a
label/value to be replaced (the entire bracketed term needs replacement).\\

\noindent Job instructions in a script start with \verb+#SBATCH+ prefix, for example:
\begin{verbatim}
    #SBATCH --mem=100M -t 600 -J <job-name> -A <slurm account>
    #SBATCH -p pg --gpus=2 --mail-type=ALL
\end{verbatim}
%
For complex compute steps within a script, use \tool{srun}. We recommend using \tool{salloc} for interactive jobs as it supports multiple steps.
However, srun can also be used to start interactive jobs (see \xs{sect:interactive-jobs}).
%
Common and required job parameters include:
\begin{verbatim}
memory (--mem), time (-t), job-name (-J), slurm project account (-A), partition (-p), 
mail type (--mail-type), ntasks (-n), CPUs per task (--cpus-per-task).
\end{verbatim}